% Capitolo: Hateville con al massimo 3 case rosse
% Da includere con \input{hateville_3_rossi} o \include{hateville_3_rossi}
% Pacchetti necessari nel preambolo del documento principale:
%   \usepackage[italian]{babel}
%   \usepackage{amsmath, amssymb}
%   \usepackage{algorithm}
%   \usepackage{algpseudocode}

\chapter{Hateville con al massimo 3 case rosse}

\newcommand{\col}{\mathrm{col}}
\newcommand{\rosso}{\mathrm{rosso}}
\newcommand{\HVBU}{\textsc{HVBU}}
\newcommand{\HVR}{\textsc{HVR}}

\section{Definizione del coefficiente}

Siano $X_n = \{1,\dots,n\}$ le case di Hateville, ciascuna con una donazione
$d(i)$ e un colore $\col(i)$. Due case adiacenti non possono essere scelte
insieme, e si possono scegliere al massimo $3$ case rosse.

Per $i \in \{0,\dots,n\}$ e $r \in \{0,\dots,3\}$ sia
\[
  S_i^r \subseteq X_i
\]
un insieme ottimo di case mutuamente compatibili (non adiacenti) scelte tra le
prime $i$ case, con $r$ case rosse ancora disponibili. Il coefficiente è
\[
  e_i^r = D(S_i^r),
\]
dove $D(\cdot)$ indica la somma delle donazioni delle case dell'insieme.

\section{Casi base}

\paragraph{Caso $i = 0$.} Per ogni $r \in \{0,\dots,3\}$ non ci sono case, quindi
\[
  S_0^r = \emptyset \;\Longrightarrow\; e_0^r = 0.
\]

\paragraph{Caso $i = 1$, $\col(1) \neq \rosso$.} La casa può essere presa senza
consumare rosse:
\[
  S_1^r = S_0^r \cup \{1\} \;\Longrightarrow\; e_1^r = e_0^r + d(1) = d(1).
\]

\paragraph{Caso $i = 1$, $\col(1) = \rosso$.}
\begin{itemize}
  \item Se $r > 0$, la casa può essere presa consumando una rossa:
  \[
    S_1^r = S_0^{r-1} \cup \{1\} \;\Longrightarrow\; e_1^r = e_0^{r-1} + d(1) = d(1).
  \]
  \item Se $r = 0$, la casa non può essere presa:
  \[
    S_1^0 = S_0^0 \;\Longrightarrow\; e_1^0 = e_0^0 = 0.
  \]
\end{itemize}

\section{Caso ricorsivo ($i > 1$)}

\subsection*{Se $\col(i) \neq \rosso$}

Si sceglie tra prendere la casa $i$ (e quindi saltare la $i-1$) oppure non
prenderla:
\begin{align*}
  e_{i-2}^r + d(i) \geq e_{i-1}^r
    &\;\Longrightarrow\; S_i^r = S_{i-2}^r \cup \{i\}
     \;\Longrightarrow\; e_i^r = e_{i-2}^r + d(i),\\
  e_{i-2}^r + d(i) < e_{i-1}^r
    &\;\Longrightarrow\; S_i^r = S_{i-1}^r
     \;\Longrightarrow\; e_i^r = e_{i-1}^r.
\end{align*}

\subsection*{Se $\col(i) = \rosso$}

\paragraph{Caso $r = 0$.} Non ci sono rosse disponibili, quindi la casa $i$ non
può essere presa:
\[
  S_i^0 = S_{i-1}^0 \;\Longrightarrow\; e_i^0 = e_{i-1}^0.
\]

\paragraph{Caso $r > 0$.} Prendere la casa $i$ consuma una rossa:
\begin{align*}
  e_{i-2}^{r-1} + d(i) \geq e_{i-1}^r
    &\;\Longrightarrow\; S_i^r = S_{i-2}^{r-1} \cup \{i\}
     \;\Longrightarrow\; e_i^r = e_{i-2}^{r-1} + d(i),\\
  e_{i-2}^{r-1} + d(i) < e_{i-1}^r
    &\;\Longrightarrow\; S_i^r = S_{i-1}^r
     \;\Longrightarrow\; e_i^r = e_{i-1}^r.
\end{align*}

\subsection*{Ricorrenza in forma compatta}

\[
  e_i^r =
  \begin{cases}
    0 & i = 0,\\[2pt]
    d(1) & i = 1,\ \col(1) \neq \rosso,\\[2pt]
    d(1) & i = 1,\ \col(1) = \rosso,\ r > 0,\\[2pt]
    0 & i = 1,\ \col(1) = \rosso,\ r = 0,\\[2pt]
    \max\{e_{i-1}^r,\ e_{i-2}^r + d(i)\} & i > 1,\ \col(i) \neq \rosso,\\[2pt]
    e_{i-1}^0 & i > 1,\ \col(i) = \rosso,\ r = 0,\\[2pt]
    \max\{e_{i-1}^r,\ e_{i-2}^{r-1} + d(i)\} & i > 1,\ \col(i) = \rosso,\ r > 0.
  \end{cases}
\]

\section{Soluzione}

La soluzione del problema è l'insieme $S_n^3$, il cui valore è $e_n^3$.

\section{Algoritmo bottom-up}

\begin{algorithm}[ht]
\caption{Calcolo della tabella $e$ (bottom-up)}
\begin{algorithmic}[1]
\Procedure{\HVBU}{$n$}
  \State \textit{// Caso base}
  \For{$r = 0$ \textbf{to} $3$}
    \State $e_0^r \gets 0$
    \If{$\col(1) \neq \rosso$}
      \State $e_1^r \gets d(1)$
    \Else
      \If{$r > 0$}
        \State $e_1^r \gets d(1)$
      \Else
        \State $e_1^r \gets 0$
      \EndIf
    \EndIf
  \EndFor
  \State \textit{// Caso ricorsivo}
  \For{$i = 2$ \textbf{to} $n$}
    \For{$r = 0$ \textbf{to} $3$}
      \If{$\col(i) \neq \rosso$}
        \State $e_i^r \gets \max\{e_{i-1}^r,\ e_{i-2}^r + d(i)\}$
      \Else
        \If{$r > 0$}
          \State $e_i^r \gets \max\{e_{i-1}^r,\ e_{i-2}^{r-1} + d(i)\}$
        \Else
          \State $e_i^r \gets e_{i-1}^r$
        \EndIf
      \EndIf
    \EndFor
  \EndFor
  \State \Return $e$
\EndProcedure
\end{algorithmic}
\end{algorithm}

\paragraph{Complessità.} La tabella ha $4(n+1)$ celle e ogni cella si calcola in
tempo costante, quindi tempo e spazio sono $\Theta(n)$.

\section{Algoritmo di ricostruzione}

La procedura va chiamata come $\HVR(n, 3, e)$ e stampa gli indici delle case
scelte in ordine crescente.

\begin{algorithm}[ht]
\caption{Ricostruzione della soluzione}
\begin{algorithmic}[1]
\Procedure{\HVR}{$i, r, e$}
  \If{$i = 0$}
    \State \Return
  \ElsIf{$i = 1$}
    \If{$\col(i) \neq \rosso$}
      \State \textbf{print}($i$)
      \State \Return
    \Else
      \If{$r > 0$}
        \State \textbf{print}($i$)
        \State \Return
      \Else
        \State \Call{\HVR}{$i-1, r, e$}
      \EndIf
    \EndIf
  \Else
    \If{$\col(i) \neq \rosso$}
      \If{$e_{i-1}^r \leq e_{i-2}^r + d(i)$}
        \State \Call{\HVR}{$i-2, r, e$}
        \State \textbf{print}($i$)
      \Else
        \State \Call{\HVR}{$i-1, r, e$}
      \EndIf
    \Else
      \If{$r > 0$}
        \If{$e_{i-2}^{r-1} + d(i) \geq e_{i-1}^r$}
          \State \Call{\HVR}{$i-2, r-1, e$}
          \State \textbf{print}($i$)
        \Else
          \State \Call{\HVR}{$i-1, r, e$}
        \EndIf
      \Else
        \State \Call{\HVR}{$i-1, r, e$}
      \EndIf
    \EndIf
  \EndIf
\EndProcedure
\end{algorithmic}
\end{algorithm}

\paragraph{Complessità.} Ogni chiamata ricorsiva riduce $i$ di almeno $1$ e fa
lavoro costante, quindi la ricostruzione costa $O(n)$.